\documentclass[11pt,letterpaper]{article}
\usepackage[margin=0.88in,headheight=15pt]{geometry}
\usepackage[T1]{fontenc}
\usepackage{lmodern,microtype,amsmath,amssymb,amsthm,mathtools}
\usepackage{booktabs,tabularx,array,xcolor,fancyhdr}
\usepackage[colorlinks=true,linkcolor=blue!45!black,urlcolor=blue!50!black,citecolor=blue!45!black]{hyperref}
\definecolor{ink}{RGB}{28,49,68}
\pagestyle{fancy}
\fancyhf{}
\setlength{\parindent}{0pt}
\setlength{\parskip}{6pt}
\setlength{\emergencystretch}{2em}
\numberwithin{equation}{section}
\theoremstyle{plain}
\newtheorem{theorem}{Theorem}[section]
\newtheorem{lemma}[theorem]{Lemma}
\newtheorem{corollary}[theorem]{Corollary}
\newtheorem{proposition}[theorem]{Proposition}
\theoremstyle{remark}
\newtheorem{remark}[theorem]{Scope}
\DeclareMathOperator{\Var}{Var}
\DeclareMathOperator{\Cov}{Cov}
\DeclareMathOperator{\supp}{supp}
\DeclareMathOperator{\Vol}{Vol}
\DeclareMathOperator{\Tr}{Tr}
\DeclareMathOperator{\diver}{div}
\newcommand{\E}{\mathbb E}
\newcommand{\R}{\mathbb R}
\newcommand{\N}{\mathbb N}
\newcommand{\dd}{\,\mathrm d}
\newcommand{\norm}[1]{\left\lVert#1\right\rVert}
\newcommand{\abs}[1]{\left|#1\right|}
\newcommand{\rref}[1]{\cite{W}, #1}
\newcommand{\OH}{\mathcal O_H}
\newcommand{\Dhat}{\widehat D}
\newcommand{\KL}{D_{\mathrm{KL}}}
\title{Nonlinear dispersion and unequal spatial weights}\author{}\date{}
\begin{document}\maketitle
\section{The original source and what an exact representation must preserve}

Fix $G=SU(N)$, $N\ge2$, a periodic isotropic four-dimensional lattice of fine spacing
$a$, and normalized product Haar measure $\lambda$. Use the conventions of PD12:
\begin{equation}
 S_0(U)=\sum_p s_p(U),\qquad s_p=N-\Re\Tr U_p,\qquad
 \dd\mu_\beta=Z_\beta^{-1}e^{-\beta S_0}\dd\lambda,
 \quad \beta=c_\beta g^{-2},\ c_\beta>0.
 \label{eq:wilson}
\end{equation}
Let $X_b^A$ be the original invariant link derivatives. Write
\[
 F_b^A=X_b^AS_0,\qquad F_{b;\mu}^A=X_b^AS_0^{(\mu\nu_b)},
 \qquad F_b^A=\sum_{\mu\ne\nu_b}F_{b;\mu}^A.
\]
With every pointed test offset and finite hypercubic average retained, define
\begin{align}
 G_H(f)&=\sum_{b,A}f_b\left\{\sum_{\mu\ne\nu_b}(F_{b;\mu}^A)^2
                       -\frac13(F_b^A)^2\right\},\label{eq:GH}\\
 \OH(f)&=g^{-2}a^{-2}G_H(f)=\frac{a^{-2}}{c_\beta}\,\beta G_H(f).
 \label{eq:normalization}
\end{align}
For a constant test, the global sum is already invariant under that finite averaging.
For a test with a plateau, every determining stencil used in the proof below lies
inside the plateau, so the same local calculation applies to the full average.

The original two-label multiplier is
\begin{equation}
 \exp_{\mathcal J}\{\eta_1O_1+\eta_2O_2\},\qquad \eta_i^2=0.
 \label{eq:source}
\end{equation}
Its bare mixed action coefficient is zero. After an exact source-free operation $T$,
the actual clean mixed coefficient, where its original body $T\rho$ is nonzero, is
\begin{equation}
 A_{12}^{\mathrm{out}}
 =\frac{T(\rho O_1O_2)}{T\rho}
  -\frac{T(\rho O_1)T(\rho O_2)}{(T\rho)^2}.
 \label{eq:sourceout}
\end{equation}
Still-live selected histories retain the undivided numerator coefficients; a
normalized positive conditional law is not inferred for an individual selected term.

For the specified decomposition $O_i=t_i+r_i$, with
$t_i\rho=-\diver(\rho V_i)$, the exact compensated coefficient is
\begin{equation}
 b_{12}=A_{12}+\frac12(V_1t_2+V_2t_1)+V_1r_2+V_2r_1.
 \label{eq:b12}
\end{equation}
Terminal integration gives
\begin{equation}
 \E A_{12}+\Cov(O_1,O_2)=\E b_{12}+\Cov(r_1,r_2).
 \label{eq:exactcov}
\end{equation}
These are PD12.13 and PD12.16--18. Complete auxiliary normalization removes its
own closed-fiber responses, not the physical right side of
\eqref{eq:exactcov}. No Ward-completed source is substituted below.

\section{A nonlinear dispersion theorem for the actual Wilson source}
\label{sec:dispersion}

The variance lower bound in PD12.40 does not alone imply non-tightness: a random
variable can have divergent variance because of very rare excursions. The argument
here controls the probability of every bounded interval directly.

\begin{theorem}[Concentration of the original $H$ source]
\label{thm:concentration}
Let $f$ be a fixed bounded real physical test equal to one on a box of positive
four-volume. Assume that this box and its fixed microscopic determining stencils
embed in the torus. There are constants $a_0,\beta_0,c_f,C_f>0$, independent of the
cutoff and of the ambient isotropic torus once those records fit, such that for
$0<a<a_0$, $\beta\ge\beta_0$, and $L\ge0$,
\begin{equation}
 \sup_{c\in\R}\mu_\beta\{\abs{\OH(f)-c}\le L\}
 \le \min\left(1,C_f\left[
          \frac{\log(2+\beta)}{\beta}+a^2+a^4L\right]\right).
 \label{eq:AC}
\end{equation}
The source, measure, and normalization are exactly
\eqref{eq:wilson}--\eqref{eq:normalization}. In particular, if
$\beta(a)\to\infty$ and $L_a=o(a^{-4})$, the left side tends to zero.
\end{theorem}

\subsection{Uniform local probability of a good bath}

For completeness, the elementary nonlinear inputs are proved here. Restrict every
link to a Haar ball of radius $c\beta^{-1/2}$ around the identity. Each plaquette
then has action at most $C\beta^{-1}$ and each ball has Haar mass at least
$c'\beta^{-\dim G/2}$. Since the plaquette/link incidence ratio is fixed,
\begin{equation}
 -\log Z_\beta\le C\#E\log(2+\beta).
 \label{eq:partition}
\end{equation}
The function $\E_\beta S_0$ decreases with $\beta$, as its derivative is
$-\Var_\beta S_0$. Therefore
\[
 \frac\beta2\E_\beta S_0
 \le \log Z_{\beta/2}-\log Z_\beta\le-\log Z_\beta.
\]
Translation and coordinate-permutation symmetry of the isotropic torus give
\begin{equation}
 \E_\beta s_p\le C\frac{\log(2+\beta)}{\beta}
 \quad\hbox{for each original plaquette }p.
 \label{eq:plaquette}
\end{equation}
For any fixed microscopic contractible patch $Q$ and fixed small $\varepsilon>0$,
\begin{equation}
 \Pr\{s_p\le\varepsilon^2\text{ for every }p\in Q\}
 \ge1-C_Q\varepsilon^{-2}\frac{\log(2+\beta)}{\beta}.
 \label{eq:localgood}
\end{equation}

For a link $e$, fix a cube containing its full $H$ determining neighborhood.
Call the conditioned bath good if a vertex gauge puts all bath links in that
neighborhood within $\delta$ of the identity. This is bath-measurable: minimize
the maximum transformed distance over the compact finite vertex gauge group.
Small plaquettes on the containing cube imply this property. In a fixed
spanning-tree gauge, every bounded-length fundamental cycle is a product of a
fixed number of conjugated plaquettes. For unitary matrices, telescoping bounds
its distance from the identity by the sum of the plaquette distances. Taking
$\varepsilon$ sufficiently small in terms of $\delta$ proves
\begin{equation}
 \Pr\{e\text{ has a bad bath}\}\le
 \varepsilon_\beta,
 \qquad \varepsilon_\beta=C_\delta\frac{\log(2+\beta)}{\beta}.
 \label{eq:badbath}
\end{equation}
The witnessing gauge can use a full configuration to prove the implication;
the good-bath event itself quantifies only over the conditioned links.

\subsection{Independent conditional links and the complete local observable}

The link interaction graph has uniformly bounded degree. Choose a deterministic
family $\mathcal F_a$ in an interior sub-box of the plateau whose members have
pairwise graph distance larger than the fixed diameter of every complete
determining stencil. A greedy packing gives
\begin{equation}
 m_a:=\#\mathcal F_a\ge c_f a^{-4}.
 \label{eq:packing}
\end{equation}
Condition on all links outside $\mathcal F_a$, denoting the bath by $b$.
No plaquette reads two free links, so the original conditional Wilson law is a
product of one-link laws. Every summand of the full $G_H(f)$ reads at most one
free link. Hence, exactly,
\begin{equation}
 G_H(f)=G_0(b)+\sum_{e\in\mathcal F_a}G_e(U_e;b).
 \label{eq:additive}
\end{equation}
All summands that vary with a selected free link have test value one. No part of
the ordinary partner or current contribution has been removed.

When its bath is exactly flat, put $s(U)=N-\Re\Tr U$. At the free link its three
plane forces are $2\nabla s(U)$ and its total force is $6\nabla s(U)$, giving
$3\cdot4-36/3=0$. There are eighteen distinct other links in its six incident
plaquettes. Each contributes $(1-1/3)|\nabla s(U)|^2$. Thus
\begin{equation}
 G_e^{\mathrm{flat}}(U)=12|\nabla s(U)|^2+\text{constant}.
 \label{eq:flatlocal}
\end{equation}
This is the original physical $H$ calculation PD12.49, not the different
total-current observable considered in~\cite[pp.~1182--1194]{W}.

In exponential coordinates $x$ near the identity, the Hessian of
$12|\nabla s|^2$ is positive definite. Finite-word continuity gives a fixed chart
and $\kappa_H>0$ such that
\begin{equation}
 D_x^2G_e(x;b)\ge\kappa_HI
 \label{eq:Hhessian}
\end{equation}
for every sufficiently small good bath. The constants are uniform in $e,a,\beta$
and the ambient torus.

On a good bath the exact conditional action, up to a constant, is
\[
 S_e(U;b)=6N-\Re\Tr(UH_b),\qquad \norm{H_b-6I}\le C\delta.
\]
At $H_b=6I$ it has a unique minimum at the identity, a positive Hessian there,
and a gap away from a fixed small chart. Compactness and the implicit function
theorem give these properties uniformly for all good baths. Let $x_*(b)$ be the
minimum. Its minimum-subtracted normalizer lies between
$c\beta^{-d/2}$ and $C\beta^{-d/2}$, $d=\dim G$, by local Gaussian integration
and the uniform exterior gap.

In $z=\sqrt\beta(x-x_*)$, choose a fixed cube $Q=[-r,r]^d$ contained in that
chart for every $\beta\ge\beta_0$. The actual rescaled conditional density obeys
\begin{equation}
 p_{e,\beta,b}(z)\ge c_0>0\quad(z\in Q).
 \label{eq:densitylower}
\end{equation}
Define $q_{e,\beta,b}(z)=\beta G_e(x_*+\beta^{-1/2}z;b)$. Then
\begin{equation}
 \partial_{z_1}^2q_{e,\beta,b}\ge\kappa_H\quad\hbox{on }Q.
 \label{eq:qhessian}
\end{equation}
The possibly large linear term in this rescaled function is retained.

\subsection{A one-link concentration deficit}

\begin{lemma}[Strongly convex bands]
\label{lem:band}
If $u''\ge\kappa>0$ on an interval and $I$ has length $h$, then
\[
 \operatorname{Leb}\{x:u(x)\in I\}\le2\sqrt{2h/\kappa}.
\]
\end{lemma}
\begin{proof}
Split at a possible minimum into at most two monotone branches. On an increasing
branch an interval of length $l$ in the band has endpoint difference at least
$\kappa l^2/2$, because its initial derivative is nonnegative. On a decreasing
branch integrate backward from its right endpoint. Each branch therefore has
length at most $\sqrt{2h/\kappa}$. The argument also covers a minimum outside
the original interval.
\end{proof}

Apply the lemma to every $z_1$ slice of $Q$. Choose $h_0>0$ small enough that the
inverse image of any interval of length $h_0$ occupies at most half the volume
of $Q$. Equation~\eqref{eq:densitylower} then yields a fixed $p_0>0$ such that
\begin{equation}
 \mathcal Q_{h_0}(\beta G_e\mid b)
 :=\sup_x\Pr\{\beta G_e\in(x,x+h_0]\mid b\}\le1-p_0
 \label{eq:deficit}
\end{equation}
for every good link. This proof uses a density lower bound and strong convexity;
a variance lower bound would be insufficient.

\subsection{Conditional summation and proof of the theorem}

Let $B_a(b)$ be the number of bad links in $\mathcal F_a$. By
\eqref{eq:badbath}, $\E B_a\le m_a\varepsilon_\beta$. Markov's inequality gives
\begin{equation}
 \Pr\{B_a>m_a/2\}\le2\varepsilon_\beta.
 \label{eq:badfraction}
\end{equation}
No independence of these bath events is used.

On the complement, at least $m_a/2$ conditionally independent variables satisfy
\eqref{eq:deficit}. The Kolmogorov--Rogozin concentration inequality, in the
equal-radius form stated in Theorem 1.1 of~\cite{kr}, gives
\[
 \mathcal Q_{h_0}\left(\sum_{e\in\mathcal F_a}\beta G_e\ \middle|\ b\right)
 \le C(p_0m_a)^{-1/2}.
\]
Additional independent bad-link summands do not increase the concentration
function, since convolution averages interval probabilities. Covering an
interval of length $h$ by at most $1+h/h_0$ intervals yields
\begin{equation}
 \mathcal Q_h\left(\sum_e\beta G_e\ \middle|\ b\right)
 \le C_f(1+h)m_a^{-1/2}.
 \label{eq:sumconcentration}
\end{equation}
By \eqref{eq:normalization}, an interval of radius $L$ in $\OH(f)$ becomes one
of length $2c_\beta a^2L$ in this sum. The bath-dependent constant in
\eqref{eq:additive} only translates the interval. Equations
\eqref{eq:packing}, \eqref{eq:badfraction}, and \eqref{eq:sumconcentration}
give \eqref{eq:AC}. This proves Theorem~\ref{thm:concentration}.

\subsection{Characteristic functions and the necessary scale of scalar rescaling}

There is also a direct Fourier conclusion, with no concentration theorem needed.

\begin{proposition}[Fourier decay away from a shrinking interval]
Under the hypotheses of Theorem~\ref{thm:concentration},
\begin{equation}
 \left|\E_\beta e^{it\OH(f)}\right|
 \le C_f\frac{\log(2+\beta)}{\beta}+e^{-c_fa^{-4}}
 \qquad(|t|\ge C_fa^2).
 \label{eq:Fourier}
\end{equation}
\end{proposition}
\begin{proof}
Equation~\eqref{eq:densitylower} minorizes each good-link law by a fixed positive
multiple $p_*\le1/2$ of uniform measure on $Q$. For fixed transverse coordinates,
the elementary second-derivative oscillatory estimate gives
\[
 \left|\int_{-r}^re^{isq(t,z')}\dd t\right|\le C(|s|\kappa_H)^{-1/2}.
\]
To verify it, split where $|q'|\le v$: that interval has length at most
$2v/\kappa_H$. On the two remaining intervals, integrate by parts, using the
monotonicity of $q'$; their combined bound is $C/(|s|v)$. Choose
$v=(\kappa_H/|s|)^{1/2}$. Averaging over the other coordinates gives the same
bound for uniform measure on $Q$, up to its fixed normalization.
For $|s|\ge S_0$, the good-link characteristic function consequently has modulus
at most $1-p_*/2$. In the physical source, $s=ta^{-2}/c_\beta$ exactly. On the
good-count event there are at least $m_a/2$ such factors. Conditional independence,
\eqref{eq:packing}, and \eqref{eq:badfraction} prove \eqref{eq:Fourier}.
\end{proof}

\begin{corollary}[No tight bare law after deterministic centering]
If $\beta(a)\to\infty$, then for every deterministic $c_a$ and every fixed $L$,
\[
 \Pr\{|\OH(f)-c_a|\le L\}\longrightarrow0.
\]
More generally, tightness of $z_a\OH(f)-c_a$, for deterministic real $z_a,c_a$,
requires $|z_a|=O(a^4)$.
\end{corollary}
\begin{proof}
The first conclusion follows from \eqref{eq:AC}. If $|z_a|/a^4$ were unbounded,
take a subsequence along which it tends to infinity and apply that estimate with
$L/|z_a|$. Every fixed interval would again have probability tending to zero,
contradicting tightness. This is a necessary condition, not a sufficiency or a
central limit theorem.
\end{proof}

\begin{remark}
These statements concern the unchanged positive Wilson random variable. They
do not rule out renormalized separated correlations, operator mixing, a local
contact-subtracted generating functional, or an Osterwalder--Schrader hierarchy
that is not the moment hierarchy of this bare random distribution. In particular,
the scalar rescaling conclusion must not be imposed on all possible renormalized
Yang--Mills composite fields. No running-coupling trajectory has been proved by
assuming $\beta(a)\to\infty$ in this consequence.
\end{remark}

\section{Strong convexity and covariance on the admitted native reference}
\label{sec:native}

This section uses the exact admitted reference of NB12.1--8, not an arbitrary
selected numerator. Its actual relative forest has $q$ free group coordinates,
$m=\dim G$, and $n=mq$. The original real domain and density are
\begin{align}
 \Omega_R&=\{B:|B_e|_*<R\text{ for every original free coordinate }e\},
       &R&=M_0p_0(h),\label{eq:domain}\\
 f(B)&=e^{-\frac12B^TQB+\ell\cdot B+W_{\rm nl}(B)}J(hB),
       &K(\dd B)&=\Dhat^{-1}\mathbf1_{\Omega_R}f(B)\dd B.
 \label{eq:nativedensity}
\end{align}
Here $h$ denotes the native current coupling to distinguish it from any saved
physical-source coupling. The genuine raw denominator is $D=P\Dhat$;
$P$ retains the classical action, $n\log h$, and original Haar factors.
The product of norm balls is convex. Its corners are allowed.

On this admitted branch, NB12 supplies $Q\ge\lambda_MI$ and, with
$\epsilon(h)=C_{M,\delta}|\Lambda|_{\rm nat}h^{1-\delta}$,
\[
 |\partial_i\partial_jW_{\rm nl}|\le c_{M,2}\epsilon(h)R^{-2},
 \qquad \|D^2\log J(hB)\|_{\rm op}\le c_{J,m}h^2.
\]
The latter bound uses the actual block-diagonal product Haar Jacobian.
The former uses NB12's enlarged native target range and summed nonlinear action
estimate; finite dimension and analyticity alone are not substitutes.

Set
\begin{equation}
 \eta_h=n c_{M,2}\epsilon(h)R^{-2}+c_{J,m}h^2,
 \qquad \kappa_0=\lambda_M-\eta_h.
 \label{eq:kappa}
\end{equation}
For sufficiently small $h$, the actual bounds $n\le n_M$ and
$|\Lambda|_{\rm nat}\le C M^4$ give $\eta_h\le\lambda_M/2$.
These counts refer to that native component; they cannot be substituted for a
growing union of components. With $U=-\log f$,
\begin{equation}
 D^2U=Q-D^2W_{\rm nl}-D^2\log J(hB)\ge\kappa_0I.
 \label{eq:convex}
\end{equation}
The full linear term remains in the measure and contributes zero to its Hessian.

\begin{theorem}[Native Poincare and covariance bounds]
\label{thm:poincare}
Under \eqref{eq:kappa} with $\kappa_0>0$, every real $F\in H^1(K)$ obeys
\begin{equation}
 \Var_KF\le\kappa_0^{-1}K|\nabla_BF|^2.
 \label{eq:poincare}
\end{equation}
Consequently
\begin{equation}
 |\Cov_K(F,G)|\le\kappa_0^{-1}
 (K|\nabla_BF|^2)^{1/2}(K|\nabla_BG|^2)^{1/2}.
 \label{eq:covnative}
\end{equation}
The domain, density, and sharp selector in these inequalities are the originals.
\end{theorem}
\begin{proof}
First use a smooth convex domain and smooth data. Put $H=F-KF$ and solve
$-Lu=H$ with $L=\Delta-\nabla U\cdot\nabla$ and $\partial_nu=0$.
The Neumann condition is imposed on the auxiliary solution, not on $F$.
Weighted integration by parts gives $KH^2=K(\nabla H\cdot\nabla u)$.
The integrated Bochner identity is
\begin{align*}
 K(Lu)^2={}&K\|D^2u\|_{\rm HS}^2
    +K\langle\nabla u,D^2U\nabla u\rangle\\
 &+\Dhat^{-1}\int_{\partial\Omega}
       \mathrm{II}(\nabla_Tu,\nabla_Tu)f\dd\sigma.
\end{align*}
The outward second fundamental form is nonnegative for the convex domain, so
$\kappa_0K|\nabla u|^2\le KH^2$. Cauchy--Schwarz gives
\eqref{eq:poincare}. Approximate the original convex norm-ball product by smooth
convex domains inside the admitted range, retaining the same density. Dominated
convergence for smooth $F$, followed by Sobolev approximation, proves the
inequality on the original sharp domain. This is a proof approximation, not a
change to the physical kernel. Covariance Cauchy--Schwarz gives
\eqref{eq:covnative}.
\end{proof}

\subsection{Original coordinate counts and moments}

Let $L_0=|\nabla U(0)|$. On the original boundary, $B\cdot n\ge0$ almost everywhere.
The divergence theorem therefore gives
\[
 n-K(B\cdot\nabla U)
 =\Dhat^{-1}\int_{\partial\Omega_R}(B\cdot n)f\dd\sigma\ge0.
\]
Since $B\cdot\nabla U(B)\ge\kappa_0|B|^2-L_0|B|$,
\begin{equation}
 K|B|^2\le M_2:=\frac{2n}{\kappa_0}+\frac{L_0^2}{\kappa_0^2}.
 \label{eq:moment2}
\end{equation}
Applying the same argument to $B|B|^{2p-2}f$ gives, for integers $p\ge1$,
\begin{equation}
 K|B|^{2p}\le
 \left(\frac2{\kappa_0}\right)^p
 \prod_{j=1}^p\left(n+2j-2+\frac{L_0^2}{2\kappa_0}\right).
 \label{eq:moments}
\end{equation}
Indeed the boundary flux remains nonnegative, and
$L_0|B|\le\kappa_0|B|^2/2+L_0^2/(2\kappa_0)$ closes the recursion.
The drift is also explicit:
\[
 L_0\le|\ell|+\sqrt n\,c_{M,1}\epsilon(h)R^{-1},
\]
because the Haar logarithm has zero first derivative at zero. There is no
$R^{2p}$ moment loss. When $q=0$ the native law is a point mass and all internal
variances vanish.

The boundary in this proof was retained with its sign. In particular, the
interior-score identity NB12.18--19 still contains its original signed boundary
partner. Poincare does not turn a physical current into an internal coordinate
derivative or remove a kernel/reference response.

\subsection{An Analysis of the native linear term}

A bound $|\ell|\le C_M$ alone does not imply
$f\asymp e^{-B^TQB/2}$ with a bounded multiplicative ratio on expanding domains:
$e^{-B^2/2+B}$ is a counterexample. The original admitted native input is stronger.
Ba\l aban's R-II, printed p.357, (1.5)--(1.6), estimates the \emph{whole linear
form on its original cutoff}; the separation exponential makes that complete
term bounded, for example by one~\cite{balabanRII}. This is the input needed for
NB12.6's same-$Q$ sandwich, and should be cited explicitly.

If only $|\ell|\le L$ and $|W_{\rm nl}+\log J|\le E$ are known, the safe
replacement is Young's inequality:
\begin{align*}
 e^{-E-L^2/\lambda}e^{-(\Lambda/2+\lambda/4)|B|^2}
 \le f(B)\le e^{E+L^2/\lambda}e^{-\lambda|B|^2/4},
 \qquad \lambda I\le Q\le\Lambda I.
\end{align*}
The Hessian and radial-moment arguments above require only the actual drift,
so they avoid that stronger comparison altogether.

\section{Actual unequal-weight response with a quadratic entropy bound}
\label{sec:tilt}

Use the genuine UW12 same-target, same-forest, same-selector comparison on its
verified top-scale branch. Its centered action difference is
\begin{equation}
 w(B)=h^{-2}\{A(\Delta\zeta,W_B)-A(\Delta\zeta,U_0)\}.
 \label{eq:tilt}
\end{equation}
To avoid confusing a response coefficient with the fine mesh $a$, write the
UW12 constants as $\alpha_w,\delta_w\ge0$:
\begin{equation}
 |w(B)|\le\alpha_w|B|+\delta_w|B|^2,
 \qquad |\nabla_Bw|\le\alpha_w+2\delta_w|B|.
 \label{eq:wgrad}
\end{equation}
The verified UW12.20h internal derivative estimate gives, with the original
target count and norm equivalence,
\begin{equation}
 H_w:=\sup_{\Omega_R}\|D^2w\|_{\rm op}
 \le C_{M,n}\{\alpha_w/R+\delta_w(1+c_*\sqrt q)\}.
 \label{eq:whess}
\end{equation}
No slower mixed-shell sum is inferred from this top-scale estimate.

Suppose $H_w\le\kappa_0/2$ and put
\[
 K_t=\frac{e^{-tw}}{K_0(e^{-tw})}K_0,\quad
 \kappa=\kappa_0-H_w,\quad L_*=L_0+\alpha_w,\quad
 M_{2,*}=2n/\kappa+L_*^2/\kappa^2.
\]
This is the exact interpolation on the original integration cycle. Each
scalar potential $\mathcal V_t=-\log f_0+tw$ has Hessian at least $\kappa I$.
Theorem~\ref{thm:poincare} and
\eqref{eq:moment2} give
\begin{equation}
 \Var_{K_t}(w)\le B_w:=\kappa^{-1}
             (2\alpha_w^2+8\delta_w^2M_{2,*}),\qquad 0\le t\le1.
 \label{eq:varw}
\end{equation}
Let $I(t)=\log K_0(e^{-tw})$. Direct differentiation gives
$I''(t)=\Var_{K_t}(w)$. Thus
\begin{align}
 \KL(K_0\Vert K_1)&=\int_0^1(1-t)I''(t)\dd t\le B_w/2,\label{eq:entropy1}\\
 \KL(K_1\Vert K_0)&=\int_0^1tI''(t)\dd t\le B_w/2.\label{eq:entropy2}
\end{align}
For an actual common carrier $F$ with square-integrable internal derivative,
\begin{equation}
 |K_1F-K_0F|\le
 \frac{\sqrt{2\alpha_w^2+8\delta_w^2M_{2,*}}}{\kappa}
 \int_0^1(K_t|\nabla_BF|^2)^{1/2}\dd t.
 \label{eq:carrier}
\end{equation}
This follows by differentiating $K_tF$ and applying
\eqref{eq:covnative}. If the carrier changes, $K_1(F_1-F_0)$ remains. A
parameter-dependent carrier also has its literal $K_t(\partial_tF_t)$ response.

The corresponding absolute denominator formula is
\begin{equation}
 \log\frac{D_1}{D_0}
 =-h^{-2}A(\Delta\zeta,U_0)-K_0w+\mathcal E,
 \qquad 0\le\mathcal E\le B_w/2.
 \label{eq:rawratio}
\end{equation}
The raw classical term and the first-order expectation remain. Equal coordinate
prefactors cancel only in this same-record ratio. Independent scalar increments
and complete numerator carriers have not been bounded by normalization.

\section{Physical \texorpdfstring{$H$}{H} and its lower rows on a complete native interior record}
\label{sec:physicalnative}

The Poincare inequality becomes useful only after the derivative of the actual
carrying observable is estimated. Here this is possible in a restricted native
sector. Keep the actual reference family $U_B=U_{C=hB}$ from PE12.19--24.
Let $D\Subset D^+$ be a complete top-scale record with a fixed positive current-unit
collar. All original current stars, test shifts, and stencils in $D^+$ must lie
inside the relaxed reference's free interior $\Lambda$. R-II (1.4) then gives
$J(U_0)=0$ on those stars. This is stronger than merely being in an admitted
regular region.

In the actual common representative, the native response and the current
conversion RG13--14/LM38--41 give, with
$u(B)=\varepsilon_{\rm bg}+h|B|_2$,
\begin{align}
 \|A_B\|_{C^1(D^+)}&\le C_Mu(B),\label{eq:nativeA}\\
 \|\partial_vA_B\|_{C^1(D^+)}
   +\|\partial_vJ(U_B)\|_{\infty,D^+}
 &\le C_Mh\|v\|_\infty.\label{eq:nativeJ}
\end{align}
Here $\varepsilon_{\rm bg}=C_MR_k\varepsilon_k$ is the original real background
bound, and the target supremum norm is the actual one. Integrating
\eqref{eq:nativeJ} from zero gives $\|J(U_B)\|_\infty\le C_Mh|B|_2$.
The input for \eqref{eq:nativeJ} is the native variational response, including
its second-order rows; it is not inferred from $C^1$ regularity alone.

\subsection{Converting total-current control to directional-current control}

On the fine lattice let $F_{\rm lin}=d_aA$. The exact Wilson plane force has
the form
\begin{equation}
 J_{\nu;\mu}(A)=\delta^*_{a,\mu}(F_{\rm lin})_{\mu\nu}
                   +N_{\nu;\mu}(A),
 \qquad J=d_a^*d_aA+N(A).
 \label{eq:planeforce}
\end{equation}
On a fixed small $C^1$ chart,
\begin{equation}
 \|N(A)\|_\infty+\sum_{\mu\ne\nu}\|N_{\nu;\mu}(A)\|_\infty
 \le C\|A\|_{C^1},\qquad
 \|DN(A)[H]\|_\infty+\sum\|DN_{\nu;\mu}(A)[H]\|_\infty
 \le C\|H\|_{C^1}.
 \label{eq:nonlinearremainder}
\end{equation}

Here is why this estimate has no hidden second derivative divided by $a$.
Write the four plaquette exponents, using signed anti-Hermitian potentials, as
$aA$, $aB+a^2b$, $-aA-a^2c$, $-aB$, where $b,c$ are the actual first differences.
Expanding the finite product and its analytic Wilson flux yields
\[
 F_a^W(A)=d_aA+[A_\mu,A_\nu]+a\mathcal R_a(A,D_aA).
\]
The linear term is exact, the order-two baseline is the commutator, and every
shift correction contains an additional $a$. On a small chart the remainder
and its first field derivative are bounded using only $C^1$ norms. Its
divergence is controlled by
$\|a\delta_a^*D\mathcal R_a[H]\|_\infty\le2\|D\mathcal R_a[H]\|_\infty$
per orientation. Differences of the commutator use only first differences,
and the exact parallel-transport corrections are bounded connection terms.
This proves the plane-resolved form of LM38--41 used in
\eqref{eq:nonlinearremainder}.

Since $d_aF_{\rm lin}=0$, the local lattice Hodge estimate gives
\begin{equation}
 \|\nabla_aF_{\rm lin}\|_{L^2_a(D)}
 \le C_{\rm col}\{\|d_a^*F_{\rm lin}\|_{L^2_a(D^+)}
                         +\|F_{\rm lin}\|_{L^2_a(D^+)}\}.
 \label{eq:hodge}
\end{equation}
To prove it, sample a smooth cutoff $\chi=1$ on the complete one-stencil
enlargement of $D$, with every shifted support inside $D^+$. This is possible
for $a$ below a fixed fraction of the original collar width. Extend $\chi F$
by zero. The exact lattice Hodge identity is
\[
 \|\nabla_a(\chi F)\|_2^2
 =\|d_a(\chi F)\|_2^2+\|d_a^*(\chi F)\|_2^2.
\]
Use the shifted discrete product rules and the bound on the first difference
of $\chi$, which depends on the original collar width. Every shifted stencil
lies in $D^+$. This proves \eqref{eq:hodge}, componentwise in a Hermitian norm.

Let $V_D=a^4N_D$, where $N_D$ is the actual number of fine bond/stencil
positions of the padded record; fixed plane and color incidences are retained
in the fixed constants. Apply \eqref{eq:hodge} to $A_B$ and $\partial_vA_B$,
using \eqref{eq:nativeA}--\eqref{eq:nonlinearremainder}. Then
\begin{align}
 \left(\sum_{\nu,\mu\ne\nu}\|J_{\nu;\mu}(U_B)\|_{L^2_a(D)}^2\right)^{1/2}
 &\le C_{M,\rm col}V_D^{1/2}u(B),\label{eq:dir1}\\
 \left(\sum_{\nu,\mu\ne\nu}\|\partial_vJ_{\nu;\mu}(U_B)\|_{L^2_a(D)}^2\right)^{1/2}
 &\le C_{M,\rm col}hV_D^{1/2}\|v\|_\infty.
 \label{eq:dir2}
\end{align}
The physical fine-mesh derivative is included. No microscopic count was
replaced by the target dimension $n$.

\subsection{Actual source gradients and the Casimir action term}

Let $\gamma$ and $\beta_s$ be the source's saved coupling factors. In this
theorem $a$ is the actual fine mesh of $U_B$ and of the source word evaluated
on that field, in current units. Keep PD12's physical currents
$J_b=a^{-3}X_bS_0$ and define
\[
 D_a(f)=a^4\sum_{b,A,\mu\ne\nu_b}f_b(J_{b;\mu}^A)^2,
 \quad P_a(f)=a^4\sum_{b,A}f_b(J_b^A)^2,
 \quad \OH(f)=\gamma^{-2}\{D_a(f)-P_a(f)/3\}.
\]
Cauchy--Schwarz in these literal sums, followed by the supremum over Euclidean
unit target directions $v$, gives
\begin{equation}
 \|\nabla_B\OH(f)\|_2
 \le C\|f\|_\infty\gamma^{-2}V_Dh\,u(B).
 \label{eq:gradH}
\end{equation}
The simultaneously controlled target sup norm satisfies $\|v\|_\infty\le1$;
no extra fine-coordinate factor is introduced.

The exact ordinary partner is
\begin{align}
 r_H(f)&=c_f+d_f+l_f,\label{eq:rowdecomp}\\
 c_f&=-\frac{C_F}{3\beta_s}\gamma^{-2}a^{-2}N\sum_pf_p,
 &d_f&=\gamma^{-2}D_a(f),\notag\\
 l_f&=\frac{C_F}{3\beta_s}\gamma^{-2}a^{-2}S_f,
 &S_f&=\sum_pf_ps_p,
 \qquad f_p=\sum_{b\in\partial p}f_b.\notag
\end{align}
The full constant $c_f$ stays in absolute masses. Centering removes it from
covariance, but does not remove the field-dependent $l_f$.

Let $V_Y=a^4\#Y_f$ use the literal completed plaquette-incidence support.
On the actual real native family,
$\|U_p-I\|\le Ca^2u(B)$ and
$\|\partial_vU_p\|\le Ca^2h\|v\|_\infty$.
The traceless group tangent gives
$|\partial_vs_p|\le Ca^4h u(B)\|v\|_\infty$.
Thus
\begin{align}
 \|\nabla_Bd_f\|_2&\le C\|f\|_\infty\gamma^{-2}V_Dh u(B),\label{eq:gradd}\\
 \|\nabla_Bl_f\|_2&\le C\|f\|_\infty\gamma^{-2}\beta_s^{-1}a^{-2}V_Yh u(B).
 \label{eq:gradl}
\end{align}
All source offsets and the four original plaquette incidences are present.

Define the explicit native factor
\begin{equation}
 \mathcal Q_{\rm nat}=\kappa_0^{-1}h^2
                (\varepsilon_{\rm bg}^2+h^2M_2),
 \qquad M_2=2n/\kappa_0+L_0^2/\kappa_0^2.
 \label{eq:Qnat}
\end{equation}
Equations \eqref{eq:covnative}, \eqref{eq:gradH}, \eqref{eq:gradd}, and
\eqref{eq:gradl} prove, for two real tests $f,k$,
\begin{align}
 |\Cov_K(\OH(f),\OH(k))|
 &\le C\mathcal Q_{\rm nat}\gamma^{-4}
                \|f\|_\infty\|k\|_\infty V_{D,f}V_{D,k},\label{eq:Hcov}\\
 |\Cov_K(d_f,d_k)|
 &\le C\mathcal Q_{\rm nat}\gamma^{-4}
                \|f\|_\infty\|k\|_\infty V_{D,f}V_{D,k},\label{eq:ddcov}\\
 |\Cov_K(d_f,l_k)|
 &\le C\mathcal Q_{\rm nat}\gamma^{-4}\beta_s^{-1}a^{-2}
                \|f\|_\infty\|k\|_\infty V_{D,f}V_{Y,k},\label{eq:dlcov}\\
 |\Cov_K(l_f,l_k)|
 &\le C\mathcal Q_{\rm nat}\gamma^{-4}\beta_s^{-2}a^{-4}
                \|f\|_\infty\|k\|_\infty V_{Y,f}V_{Y,k}.\label{eq:llcov}
\end{align}
Distinct saved couplings use their respective factors on the two sides.
Sources on distinct original graphs or meshes require the response and
complete-collar hypotheses to be verified separately on those graphs before
combining their gradient bounds. Substituting another mesh into the displayed
common-mesh formulas does not cover an inherited reconstructed source.

In the checked same-scale specialization $\gamma=h$, $\beta_s=c_\beta h^{-2}$,
\begin{align}
 \Var_K\OH(f)&\le C\kappa_0^{-1}\|f\|_\infty^2V_D^2
                      [ (\varepsilon_{\rm bg}/h)^2+M_2],\label{eq:samescaleH}\\
 \Var_Kr_H(f)&\le C\kappa_0^{-1}\|f\|_\infty^2
       [V_D+c_\beta^{-1}h^2a^{-2}V_Y]^2
                      [(\varepsilon_{\rm bg}/h)^2+M_2].\label{eq:samescaler}
\end{align}
These estimates separate the physical dimension-six observable from its
dimension-four lower row. They do not assert that the latter variance is
bounded independently of the original mesh.

\begin{remark}
This is a consequence of the stated native response on a complete top-scale
free-interior record. It does not identify a straddling inherited coefficient
with that record, admit simultaneous mixed localization walls, or supply its
old saved grade. In particular, the Wilson integration-by-parts identity
\eqref{eq:exactcov} cannot be copied to the native conditional $K$ without its
actual effective score and boundary/kernel companions. A saved
field-independent source factor $\Gamma$ multiplies its gradient bound by
$|\Gamma|$, its variance bound by $|\Gamma|^2$, and a covariance with a second
saved factor by the product of their absolute values. A live factor also
contributes its actual derivative.
\end{remark}

\section{A relative-weight comparison that includes all native shells}
\label{sec:flat}

There is one sector in which the unresolved small-exponent native metric sum
can be avoided. Keep a single admitted native record, its actual target map
$W_B$, graph, pins, forest, cutoff, and Haar density. Assume its relaxed
minimum $U_0$ is plaquette-flat. Let the original baseline weight $\zeta_0\ge0$
have the admitted coercive native expansion. Consider weights satisfying
\begin{equation}
 |\zeta_1(p)-\zeta_0(p)|\le\delta\zeta_0(p),\qquad0\le\delta<1.
 \label{eq:relativeweight}
\end{equation}
In particular no new support is introduced where $\zeta_0=0$.
Define $E_0(B)=h^{-2}A(\zeta_0,W_B)$ and
$w(B)=h^{-2}A(\zeta_1-\zeta_0,W_B)$. Flatness removes the classical subtraction
exactly, and positivity of each original real Wilson plaquette action gives
\begin{equation}
 |w(B)|\le\delta E_0(B).
 \label{eq:energyrelative}
\end{equation}
Every original assigned-scale shell is already included in this exact sum.

Since the derivative of each plaquette action vanishes at a flat plaquette,
the native linear action term is zero. The summed native expansion gives
\begin{equation}
 E_0(B)=\tfrac12B^TQ_0B-W_{\rm nl}(B),\qquad
 0\le E_0(B)\le\tfrac12\Lambda_M|B|^2+\epsilon_{\rm nl},
 \quad\epsilon_{\rm nl}\le1.
 \label{eq:flatE}
\end{equation}
This is the original summed native action bound, not an inference from a
finite number of target coordinates.

For fixed $p\ge0$, $0\le\theta\le\lambda_M/8$, and sufficiently small
$\delta$ depending on the actual native constants,
\begin{equation}
 K_0\left[(1+|B|)^pe^{\theta|B|^2}|e^{-w}-1|\right]
 \le C_{M,n,p,\theta}\delta.
 \label{eq:allflat}
\end{equation}
Indeed $|e^{-w}-1|\le\delta E_0e^{\delta E_0}$.
On this flat branch the original upper density bound is
$f_0\le e^{C_M}e^{-\lambda_M|B|^2/2}$ and its extracted denominator is bounded
below by integration over its fixed inner ball, retaining its actual
$n$-dimensional volume. Using \eqref{eq:flatE} and
$\delta\le\lambda_M/(4\Lambda_M)$ leaves an integrable majorant proportional
to $(1+|B|)^p(1+|B|^2)e^{-\lambda_M|B|^2/4}$.
Extending this positive majorant to $\R^n$ bounds the integral; the field map
itself is never extended beyond the original cutoff.

Let $r=C_{M,n,0,0}\delta\le1/2$. The exact same-record identities imply
\begin{equation}
 \left|\log\frac{D_1}{D_0}\right|\le2r,
 \qquad \|K_1-K_0\|_{\rm TV}\le2r,
 \label{eq:allflatratio}
\end{equation}
where total variation is $\sup_A|K_1(A)-K_0(A)|$.
The weighted form follows by writing
$dK_1/dK_0=e^{-w}/K_0(e^{-w})$.
All coordinate factors cancel only because the actual records coincide.
The raw individual denominator masses and independent scalar factors remain.
For $n=0$ the reference is the fixed zero target and this comparison is exact,
with zero error.

The same proof allows complex relative weight perturbations on the unchanged
real $B$ cycle: \eqref{eq:energyrelative} is a modulus estimate, and the denominator
ratio stays in $|z-1|\le r<1$. Hence it has a holomorphic logarithm on this
relative-weight neighborhood. For a holomorphic parameter $z$ with
$|\Delta\zeta_z|\le\delta\zeta_0$ on $|z|<r_z$, its logarithmic derivatives
at zero obey $m!C\delta r_z^{-m}$. On $|z|\le\rho r_z$, $0<\rho<1$, the bound
is $m!C\delta[(1-\rho)r_z]^{-m}$.
These are weight derivatives, not unproved internal-current lifts.

\begin{remark}
Equation~\eqref{eq:allflat} genuinely includes lower native shells in this
fixed flat-reference sector. It has no physical-distance decay. A uniform
$O(\delta)$ error per microscopic vacuum birth would still have a divergent
four-dimensional root sum. Different target maps, forests, selectors, or
complete carriers require their own comparison. Therefore this estimate
does not prove S94.
\end{remark}

\section{Why the remaining global obligations are still substantive}

\subsection{Actual mixed-history admission}

The protected-row identity of~\cite[pp.~1174--1180]{W} is exact for coefficients factoring
through the complete retained rows outside the active domain. It does not
cover a coefficient whose current star, shifted test, prefix, boundary frame,
or old determining record is touched by the next join. The native estimates
in Sections~\ref{sec:native}--\ref{sec:physicalnative} are not estimates for
that joined numerator. The unresolved target remains the actual
selector-supported prepared map in the old graded domain, with all its
simultaneous complex walls and source readouts. A generic regularity assertion
cannot replace that identification.

\subsection{All-history marked transfer}

The ordinary conditional covariance now has a quantitative mechanism on an
identified native reference, and the original $H$ function has a derivative
bound on a complete native interior record. The MT36 target nevertheless
still requires, schematically,
\begin{equation}
 \left\|P_6\mathcal L_{V_k}^{(2)}
       (R_j\circ\Phi_{j,k},V_k)\right\|
 \le Cg_kL^{-s(k-j)}(g_k/g_j)^q\|R_j\|_{\rm graded}
 \label{eq:MT}
\end{equation}
for every actual inherited stable/live coefficient, with its lower and
contact rows filed on the same support. Nothing above proves the complete
physical-current lift, inherited multiplier gradients, effective-score and
live-reference responses, or common complex chart needed for this statement.
An old inverse mesh power cannot be absorbed in a native constant or in a
tail depending only on the current coupling.

\subsection{S94 is an actual finite-size pressure comparison}

Write $p_{a,T}=\Vol(T)^{-1}\log Z_{a,T}$ for the original Wilson partition
function at one common bare coupling. If the source-zero stability already
used in the manuscript is available,
$|\log Z_{a,T}-V_{a,T}|\le C_{\rm st}\Vol(T)$, then
\begin{equation}
 \left|\left(\frac{V_{a,T_*}}{\Vol(T_*)}
             -\frac{V_{a,qT_*}}{\Vol(qT_*)}\right)
       -(p_{a,T_*}-p_{a,qT_*})\right|\le2C_{\rm st}.
 \label{eq:pressure}
\end{equation}
Thus the unproved one-sided S94 bound is equivalent, up to this stated
constant, to $p_{a,T_*}-p_{a,qT_*}\le C$. It is not a normalization convention.

On the original common Haar variable space, the exact gluing identity is
\begin{equation}
 \frac{Z_{a,qT_*}}{Z_{a,T_*}^{q^4}}
 =\E_{\mu_{a,T_*}^{\otimes q^4}}e^{-\beta_a\Delta S}.
 \label{eq:gluing}
\end{equation}
Let $t_*$ be the physical side length of the reference torus $T_*$.
For $n=t_*/a$, there are $q^4(12n^3-6n^2)$ potentially rewired plaquettes.
The direct bounded-action estimate consequently loses
$O(\beta_a/(t_*a^3))$ per physical volume. Reflection positivity gives the
opposite side, $p_{a,T_*}-p_{a,qT_*}\ge0$, not the desired uniform upper bound.
Subsequential thermodynamic compactness does not supply this estimate or the
common source normalizer used in S95--99 to transfer the noncollapse bound.

The native metric problem also remains precise. NM12 distinguishes the
available sum with decay rate $\delta_0/2$ from the actual nonlinear derivative
with rate $\delta_0/8$. A faster-decay summability estimate does not imply the
slower one. The flat relative-energy estimate in Section~\ref{sec:flat} avoids
that issue only on its stated record, and supplies no missing distance gain
for wrapped or unmatched records.

\subsection{Continuum field hierarchy and Yang--Mills identification}

Even closure of the preceding interfaces would have to be followed through
the manuscript's weak-trajectory, terminal, operator-normalization, common
distribution-bound, reflection/covariance, nondegeneracy, and Yang--Mills
identification obligations. Neither a nonzero scalar covariance nor the
finite-dimensional native results are substitutes for that hierarchy.
The dispersion theorem clarifies why the bare coincident source cannot
itself be used as a tight classical random field after only subtracting its
mean. Any contact or operator mixing prescription must be displayed and
proved compatible with the resulting hierarchy.

\section{Status}

\begin{center}
\small
\begin{tabularx}{\textwidth}{@{}p{0.24\textwidth}XX@{}}
\toprule
Obligation & Result here & Remaining requirement\\
\midrule
Original physical source & Nonlinear concentration and Fourier bounds;
centering alone cannot give a tight bare $H$ law & If this $H$ sector is retained:
explicit contact and mixing prescription for its continuum hierarchy\\[4pt]
Native ordinary covariance & Strong convexity, sharp-cutoff Poincare,
dimension-explicit moments & Actual inherited multiplier gradients and
physical-current/reference lifts\\[4pt]
Physical $H$ on native interior & Directional-current $L^2$ estimate and
covariance bound with explicit Casimir lower rows & Touched/mixed records,
saved grades and simultaneous-wall admission\\[4pt]
Native spatial-weight comparison & Quadratic entropy on the verified
top-scale branch; all-shell flat relative-weight sector & General mixed-shell
distance gain, unmatched carriers and absolute scalar comparison\\[4pt]
AE8/LM47, MT36, S94 & Restricted advances above & All three general
statements remain unproved\\[4pt]
Four-dimensional existence & No complete continuum construction obtained &
Common interacting field hierarchy and Yang--Mills identification\\
\bottomrule
\end{tabularx}
\end{center}

The exact one-link count and normalization agree with PD12
of~\cite[pp.~1182--1194]{W}. Non-tightness is proved through both
the concentration-function argument and the characteristic-function argument;
neither uses variance divergence as a proxy for non-tightness. The native
covariance argument uses the original sharp boundary and its sign,
the actual linear drift, and the separate saved source and native couplings.
The Casimir constant and field-dependent action row are distinguished
throughout. The original R-II linear-form estimate (1.6) is
the justification for the admitted same-$Q$ envelope.

 The restricted native statements explicitly inherit the
admitted native response estimates named in their hypotheses. No claim is
made that the remaining statements of~\cite{W} are established here. No priority claim is made for the auxiliary inequalities.

\newpage
\begin{thebibliography}{9}
\raggedright
\bibitem{W}
Wilson-source estimates, volumes 1--10 (\texttt{lattice\_theory/source\_estimates} in this collection), cited by source page.
Main inputs used here: PD12.1--63 (pp.~1182--1194), NB12.1--24 (pp.~1200--1204),
PE12.19--24 (pp.~1179--1180), UW12.1--28 (pp.~1205--1213), NM12.1--17 (pp.~1216--1218),
LM38--41 (pp.~780--781), RG13--14 (p.~748), S94 (pp.~618--621) and SV3--16 (pp.~630--633).

\bibitem{kr}
T. Ju\v{s}kevi\v{c}ius.
\emph{The sharp form of the Kolmogorov--Rogozin inequality and a conjecture
of Leader--Radcliffe}. arXiv:2201.09861 (2022), Theorem 1.1.
\url{https://arxiv.org/abs/2201.09861}.
Only the stated classical equal-radius concentration inequality is imported;
the Wilson application is proved in Section~\ref{sec:dispersion}.

\bibitem{balabanRII}
T. Ba\l aban.
\emph{Large field renormalization. II. Localization, exponentiation, and
bounds for the R operation}. Communications in Mathematical Physics
122 (1989), 355--392.
\url{https://doi.org/10.1007/BF01238433}.
In particular, printed pp.356--358, (1.2)--(1.9), including the complete
linear-form estimate (1.6). The paper's stated stability theorem is not
imported as a full interacting existence theorem.

\bibitem{balabanvar}
T. Ba\l aban.
\emph{The variational problem and background fields in renormalization
group method for lattice gauge theories}. Communications in Mathematical
Physics 102 (1985), 277--309.
\url{https://doi.org/10.1007/BF01229381}.
The native response inputs from Proposition 9 and (173)--(190) are used
in the scopes recorded in~\cite[pp.~1214--1218]{W}; their mixed-history extensions
are not asserted here.

\end{thebibliography}
\end{document}
